\documentclass[12pt]{article}
\usepackage[utf8]{inputenc}
\usepackage[T1]{fontenc}
\usepackage{hyperref}
\title{Kurtokha Language Resources for K-12 Curriculum Development}
\author{Compiled by Claude Sonnet 5.5 (claude-sonnet-5-5) for Andr\'as Kornai}
\date{}
\begin{document}
\maketitle

\section{Introduction}
Kurtokha (ISO 639-3 xkz) is better known as Kurt\"op (Kurtoep). It is an East Bodish language of the Sino-Tibetan family, spoken mainly in Lhuntse District, northeastern Bhutan. Speaker estimates range widely, from about 2,200 to 15,000, depending on the source. The language is better documented than its summary suggests: it has a reference grammar, a documentation archive and a recent trilingual dictionary. I found no school materials.

\section{Primary Educational Resources}
None found. I found no primer or textbook, and no school programme for Kurt\"op.

\section{Linguistic and Reference Materials}
\begin{itemize}
\item Gwendolyn Hyslop, \emph{A Grammar of Kurt\"op}: a 2011 PhD dissertation and a 2017 Brill book, described as the first descriptive grammar of Kurt\"op and the only reference grammar of an East Bodish language. \url{https://search.ub.tum.de/vufind/Record/DE-604.BV044677704}
\item Hyslop's 2011 dissertation, at ProQuest. \url{https://www.proquest.com/docview/868530580}
\item A Bhutan Broadcasting Service report of June 2025 on a 319-page Kurt\"op--English--Dzongkha dictionary by Hyslop, Kuenga Lhendup and Pema Chophyel. I did not open the dictionary itself. \url{https://www.bbs.bt/?p=223819}
\item Culture in Crisis project page, ``Documenting Kurtoep in a historical perspective: a grammar, dictionary and texts''. \url{https://cultureincrisis.org/projects/documenting-kurtoep-in-a-historical-perspective-a-grammar-dictionary-and-texts}
\item Wikipedia, ``Kurt\"op language''. \url{https://en.wikipedia.org/wiki/Kurt%C3%B6p_language}
\end{itemize}

\section{Religious and Community Resources}
\begin{itemize}
\item Joshua Project profile of Kurtokha in Bhutan. \url{https://joshuaproject.net/people_groups/12878}
\end{itemize}

\section{Digital and Audio Resources}
\begin{itemize}
\item ELDP archive ``Documentation of Kurt\"op'': audio and video recordings made between 2005 and 2009. \url{https://eldp.access.preservica.com/dk0179/}
\item Endangered Languages Project entry for Kurt\"op. \url{https://endangeredlanguages.com/node/96917}
\end{itemize}

\section{Curriculum Development Context}
The existing grammar, dictionary and archived recordings are a real base for mother-tongue teaching materials. The dictionary uses Dzongkha script and the Latin alphabet, which suggests where an orthography decision might start. Whether the language is taught in Bhutanese schools, and the details of any community orthography, were not established in my searches. These would need confirming with Hyslop's team and local speakers.
\end{document}